\chapter{Introducción y antecedentes}
\label{ch:introduction}
\section{Contexto y motivación}
Explica el problema en su contexto. Identifica a quién afecta, qué se puede hacer
actualmente y qué dificultad motiva la investigación. Apoya las afirmaciones
con fuentes adecuadas, por ejemplo \cite{example}.

Avanza del contexto general a una laguna concreta de conocimiento. Distingue
lo que demuestran estudios anteriores de lo que investigará tu tesis y explica
por qué esa laguna es relevante en la teoría o en la práctica.
\section{Pregunta de investigación y objetivos}
\textbf{Pregunta de investigación:} ¿Cómo afecta un procedimiento revisado al tiempo
de realización y a la precisión frente al proceso habitual?
\begin{itemize}
\item Establecer una referencia de comparación y criterios de evaluación explícitos.
\item Desarrollar o seleccionar un enfoque que aborde la laguna identificada.
\item Evaluar los resultados y discutir los límites de las pruebas.
\end{itemize}
\section{Trabajos relacionados}
Organiza la bibliografía por temas o enfoques. Compara sus supuestos, métodos
y hallazgos. Explica cómo tu estudio se deriva de esa comparación.
\section{Alcance y estructura de la tesis}
Delimita la población, el contexto, los materiales y los resultados del estudio.
El capítulo~\ref{ch:methods} describe el método; el capítulo~\ref{ch:findings} presenta
y discute los hallazgos. El apéndice~\ref{app:materials} contiene material complementario.
